\documentclass[12pt,a4paper]{article}

\usepackage{amsmath,amssymb,amsthm}
\usepackage{geometry}
\usepackage{hyperref}
\usepackage{booktabs}
\usepackage{url}
\usepackage{cite}
\usepackage{enumitem}
\geometry{margin=2.5cm}

\newtheorem{theorem}{Theorem}
\newtheorem{definition}{Definition}
\newtheorem{remark}{Remark}

\title{%
  \textbf{Biological Transitions as Multi-Agent Realisations of the}\\
  \textbf{Generative Operator Pipeline in Topographical Orthogonal}\\
  \textbf{Generative Theory (TO/TOGT)}\\[0.5em]
  \large\emph{Version 2 --- LaTeX source, complete references, and
  companion cross-reference added}%
}

\author{%
  Pablo Nogueira Grossi\\
  G6 LLC, Newark, New Jersey, USA\\
  ORCID: 0009-0000-6496-2186\\
  \texttt{pablogrossi@hotmail.com}\\[0.3em]
  \small Concept DOI (all versions):
    \href{https://doi.org/10.5281/zenodo.19208014}{10.5281/zenodo.19208014}\\
  \small V1 DOI:
    \href{https://doi.org/10.5281/zenodo.19208015}{10.5281/zenodo.19208015}\\
  \small Companion (fruit-fly toy model):
    \href{https://doi.org/10.5281/zenodo.19210136}{10.5281/zenodo.19210136}
}

\date{March 2026 (v1); May 2026 (v2 --- this revision)}

\begin{document}
\maketitle

\begin{abstract}
We model collective biological transitions --- including neural
oscillations, HPA-axis stress response, circadian regulation, immune
adaptation, and protein conformational change --- as multi-agent
realisations of the generative operator pipeline
$G = U \circ F \circ K \circ C$ introduced in Topographical Orthogonal
Generative Theory (TO/TOGT). Each agent undergoes local compression,
curvature intensification, loss of injectivity, and stabilisation.
The collective fixed points, contraction properties, and saturated
pitchfork bifurcation behaviour follow directly from the fixed-point
theory, contraction-mapping results, and algebraic separation theorem
proved in the author's earlier work \cite{Grossi2026_Vol1}
(HAL hal-05555216v1). The specific dm$^3$ numerical invariants
$(T^{*}=2\pi,\;\mu_{\max}=-2,\;\tau=2)$ are convenient modelling
choices matching observed biological scales; they are not derived
theorems of the operator algebra. No numerical simulations or
unverified biological claims are made.

This article is part of the Principia Orthogona / Generative Contact
Mechanics series and demonstrates a concrete multi-agent biological
application of the TO/TOGT operator pipeline. A companion paper
\cite{Grossi2026_FruitFly} specialises the framework to the
\textit{Drosophila} connectome, adding Python simulation code, Lean~4
proofs, and four figures.
\end{abstract}

\noindent\textbf{Keywords:} generative operator pipeline; TO/TOGT;
multi-agent systems; biological transitions; dynamical systems;
contraction mapping; saturated pitchfork bifurcation; connectomics;
HPA axis; generative transitions

\noindent\textbf{MSC codes:} 37C25, 37D10, 37G10, 47H10, 92B20, 92C20

% ------------------------------------------------------------------
\section{Relation to Prior Work}
\label{sec:prior}

This article applies the mathematical results established in the
author's earlier preprint \textit{Principia Orthogona, Volume One:
The Mathematics of Generative Transitions}
\cite{Grossi2026_Vol1} (HAL hal-05555216v1). There the composite
operator pipeline
\[
  G = U \circ F \circ K \circ C
\]
was introduced together with complete proofs of fixed-point existence
and uniqueness, global contraction, saturated pitchfork bifurcation,
and the algebraic trace-bound/separation theorem. The geometric
realisation of the operator pipeline on contact manifolds was developed
in \textit{Principia Orthogona, Volume Two: Contact Realisation of
Generative Transitions} \cite{Grossi2026_Vol2}
(HAL hal-05559997v1 and hal-05565024v1, both deposited 24 March 2026).

The present paper demonstrates a multi-agent biological application
of the same operator algebra. It is directly related to the following
Zenodo deposits by the author:
\begin{itemize}[leftmargin=2em, itemsep=1pt]
  \item \cite{Grossi2026_Crystal}:
    Zenodo \href{https://doi.org/10.5281/zenodo.19162013}{19162013}
    (The G6 Crystal --- architectural application)
  \item \cite{Grossi2026_DM3b}:
    Zenodo \href{https://doi.org/10.5281/zenodo.19122168}{19122168}
    (dm$^3$ operator, explicit toy model)
  \item \cite{Grossi2026_DM3}:
    Zenodo \href{https://doi.org/10.5281/zenodo.19117400}{19117400}
    (dm$^3$ operator, global dynamical analysis)
\end{itemize}
All mathematical claims in this paper follow strictly from previously
proved results. The companion paper \cite{Grossi2026_FruitFly}
specialises to the fruit-fly connectome and adds code and formal
verification.

% ------------------------------------------------------------------
\section{Operator Pipeline in Biological Systems}
\label{sec:pipeline}

\begin{definition}[Biological Agent Dynamics {\cite{Grossi2026_Vol1}}]
\label{def:agent}
Let each biological agent (neuron, endocrine cell, immune cell,
protein domain) follow a trajectory $x_{i}(t)$ in a Riemannian
manifold $(M,g)$. The collective biological state is the tuple
$(x_{1},\ldots,x_{N})$. Each agent undergoes a generative transition
governed by the composite operator
\[
  G = U \circ F \circ K \circ C,
\]
where:
\begin{itemize}[leftmargin=2em, itemsep=2pt]
  \item $C$: local orthogonal compression (neighbourhood averaging);
  \item $K$: curvature intensification (alignment and clipping);
  \item $F$: loss of injectivity (folding of local information);
  \item $U$: stabilisation (global coherence).
\end{itemize}
\end{definition}

The contraction-mapping theorem of \cite{Grossi2026_Vol1} ensures
that for sufficiently small coupling between agents, the collective
system converges to a unique fixed point. Formal Lean~4/Mathlib4
verification of the arithmetic spine of this result is in the AXLE
repository \cite{AXLE2026}, file \texttt{MultiOrbitBioSwarm.lean}.

% ------------------------------------------------------------------
\section{Fruit Fly Connectome as Multi-Orbit System}
\label{sec:connectome}

The FlyWire connectome \cite{Dorkenwald2024} provides a concrete
multi-agent biological system. Each neuron trajectory is modelled as
an identity orbit $\mathcal{O}_{i}$. The full connectome is the
multi-orbit system $S = \{\mathcal{O}_{1},\ldots,\mathcal{O}_{N}\}$.
The modular community structure is documented in \cite{Schlegel2024}.

The contraction results of \cite{Grossi2026_Vol1} imply that under
small inter-neuron coupling, the system converges to a collective
fixed point preserving local connectivity while achieving global
coherence. This is a direct biological instance of multi-orbit
identity stabilisation.

A specialised treatment of this setting --- with Python simulation
code generating four figures and Lean~4 proofs of the six-iterate
convergence bound --- appears in the companion paper
\cite{Grossi2026_FruitFly}.

% ------------------------------------------------------------------
\section{HPA-Axis Stress Response as Saturated Pitchfork}
\label{sec:hpa}

The hypothalamic-pituitary-adrenal (HPA) axis exhibits threshold
behaviour under acute stress \cite{Bhaskaran2022}. This is modelled
by the saturated pitchfork bifurcation proved in
\cite{Grossi2026_Vol1}:
\[
  x' = \lambda x - c x^{3},
\]
with clipping enforcing a bounded response. At critical coupling
$|\alpha|=0.5$, the system transitions between low-stress and
high-stress attractors. The clipping function $\kappa$ prevents
runaway amplification and ensures recovery. Pitchfork bifurcations
in coupled neural oscillators are discussed in
\cite{PerezCervera2019}.

% ------------------------------------------------------------------
\section{Discussion}
\label{sec:discussion}

The TO/TOGT operator pipeline provides a unified dynamical-systems
language for describing collective biological transitions. All
mathematical claims in this paper follow directly from rigorously
proved theorems in \cite{Grossi2026_Vol1} and the contact-geometric
realisation in \cite{Grossi2026_Vol2}. The specific dm$^3$ numerical
invariants are modelling choices matching biological scales; they are
not derived from the operator algebra.

Future work will focus on explicit multi-agent simulations and
experimental falsification of predicted bifurcation thresholds.
The companion paper \cite{Grossi2026_FruitFly} already provides
the simulation code and Lean~4 verification skeleton for the
fruit-fly setting; the same infrastructure applies to the broader
biological scope treated here.

% ------------------------------------------------------------------
\section*{Acknowledgements}

The author thanks the Lean~4/Mathlib4 community for the formal
verification infrastructure and the FlyWire consortium
\cite{Dorkenwald2024} for making connectome data publicly available.
All formal verification code is maintained in the AXLE repository
\cite{AXLE2026} at \url{https://github.com/TOTOGT/AXLE}.

% ------------------------------------------------------------------
\begin{thebibliography}{99}

\bibitem{Grossi2026_Vol1}
P.~Nogueira~Grossi,
\textit{Principia Orthogona, Volume~One: The Mathematics of Generative
Transitions}.
G6~LLC. HAL: hal-05555216v1. Zenodo:
\href{https://doi.org/10.5281/zenodo.19117400}{10.5281/zenodo.19117400}
(2026).

\bibitem{Grossi2026_Vol2}
P.~Nogueira~Grossi,
\textit{Principia Orthogona, Volume~Two: Contact Realization of
Generative Transitions}.
G6~LLC. HAL: hal-05559997v1 (2026).

\bibitem{Grossi2026_FruitFly}
P.~Nogueira~Grossi,
``Biological Transitions as Multi-Agent Realisations of the Generative
Operator Pipeline in TO/TOGT: A Fruit-Fly Connectome Toy Model.''
G6~LLC. Zenodo (concept DOI):
\href{https://doi.org/10.5281/zenodo.19210136}{10.5281/zenodo.19210136}
(2026).

\bibitem{Grossi2026_DNLS}
P.~Nogueira~Grossi,
``Differential Nonlinear Robustness of Critical States in Fibonacci and
Tribonacci Substitution Chains.''
G6~LLC. Zenodo (concept DOI):
\href{https://doi.org/10.5281/zenodo.20026942}{10.5281/zenodo.20026942}
(2026).

\bibitem{Grossi2026_Crystal}
P.~Nogueira~Grossi,
``The G6 Crystal: A dm$^3$-Derived Architectural Form for
Resonance-Stable Tall Structures.''
G6~LLC. Zenodo:
\href{https://doi.org/10.5281/zenodo.19162013}{10.5281/zenodo.19162013}
(2026).

\bibitem{Grossi2026_DM3}
P.~Nogueira~Grossi,
``The dm$^3$ Operator: Explicit Toy Model and Global Dynamical
Analysis.''
G6~LLC. Zenodo:
\href{https://doi.org/10.5281/zenodo.19117400}{10.5281/zenodo.19117400}
(2026).

\bibitem{Grossi2026_DM3b}
P.~Nogueira~Grossi,
``The dm$^3$ Operator (extended).''
G6~LLC. Zenodo:
\href{https://doi.org/10.5281/zenodo.19122168}{10.5281/zenodo.19122168}
(2026).

\bibitem{AXLE2026}
P.~Nogueira~Grossi,
\textit{AXLE: Lean~4 Formal Verification Engine for TO/TOGT}.
GitHub: \url{https://github.com/TOTOGT/AXLE} (2026).

\bibitem{Dorkenwald2024}
S.~Dorkenwald \textit{et~al.}\ (FlyWire Consortium),
``Neuronal wiring diagram of an adult brain,''
\textit{Nature} \textbf{634}, 124--138 (2024).

\bibitem{Schlegel2024}
P.~Schlegel \textit{et~al.},
``Whole-brain annotation and multi-connectome cell typing of
\textit{Drosophila},''
\textit{Nature} \textbf{634}, 139--152 (2024).

\bibitem{Winding2023}
M.~Winding \textit{et~al.},
``The connectome of an insect brain,''
\textit{Science} \textbf{379}, eadd9330 (2023).

\bibitem{Bhaskaran2022}
S.~Bhaskaran \textit{et~al.},
``Bifurcation structure of the HPA axis and implications for
stress-response dynamics,''
\textit{Journal of Theoretical Biology} \textbf{540}, 111077 (2022).

\bibitem{PerezCervera2019}
A.~P\'{e}rez-Cervera, T.~Seara, and G.~Huguet,
``The uncoupling limit of identical Hopf bifurcations with an
application to perceptual bistability,''
\textit{Journal of Mathematical Neuroscience} \textbf{9}, 7 (2019).

\bibitem{Banach1922}
S.~Banach,
``Sur les op\'{e}rations dans les ensembles abstraits et leur
application aux \'{e}quations int\'{e}grales,''
\textit{Fundamenta Mathematicae} \textbf{3}, 133--181 (1922).

\end{thebibliography}

\end{document}
